\begin{lemma}\label{lem:nonic31-known-binding}
The five degree-$9$ places surviving the finite class filter on
$Y=X_1(31)_{\mathbb F_3}$ are reductions of degree-$9$ points on
$X_1(31)_{\mathbb Q}$ whose underlying elliptic curves have torsion
$\mathbb Z/31\mathbb Z$ over their residue fields.
\end{lemma}
\begin{proof}
Let $K=\mathbb Q(a)$, where
\[
 a^9-2a^8+a^6-a^5+14a^4-28a^3+19a^2-2a-1=0,
\]
and put
\[
 x=a,\qquad
 y=\frac{9a^8-28a^7+16a^6+46a^5-11a^4+125a^3-440a^2+386a+150}{119}.
\]
This is the degree-$9$ point in van Hoeij's table with $N=31$,
source line $68$. Set
\[
 r=\frac{x^2y-xy+y-1}{x^2y-x},\qquad
 s=\frac{xy-y+1}{xy},\qquad b=rs(r-1),\qquad c=s(r-1).
\]
The curve $E:y^2+(1-c)xy-by=x^3-bx^2$ is nonsingular, and
$P=(0,0)$ has exact order $31$. The defining polynomial is irreducible
modulo $3$. Its simple roots $28$ modulo $53$ and $39$ modulo $83$
give good degree-one reductions with
\[
 \#E(\mathbb F_{53})=62,\qquad \#E(\mathbb F_{83})=93.
\]
Reduction is injective on torsion at these unramified odd primes.
Hence $\#E(K)_{\mathrm{tors}}$ divides $\gcd(62,93)=31$, and thus
$E(K)_{\mathrm{tors}}\simeq\mathbb Z/31\mathbb Z$.

We reduce $(E,kP)$ modulo $3$, for $1\leq k\leq15$, and compare the
Frobenius orbits of its Tate normal coefficients with the five surviving
places. The exact bindings are
\[
\begin{array}{c|c}
\text{factor-base index}&k\\ \hline
1878&1,5,6\\
2248&4,7,11\\
3105&3,13,15\\
3225&8,9,14\\
3323&2,10,12
\end{array}
\]
All five orbits have length $9$. The finite-field embeddings and the
complete coefficient comparisons are recorded in \texttt{binding31.json}.
The program \texttt{verify\_known31.py} checks the curve equations,
orders, both point counts, and all fifteen diamond reductions directly.
Every global diamond translate has the same underlying elliptic curve
and residue field $K$. This proves the claim.
\end{proof}

\begin{lemma}\label{lem:nonic31-rigidity}
Let $D$ be a rational effective divisor of degree $9$ on $X_1(31)$.
If its reduction modulo $3$ is one of the five places in
Lemma~\ref{lem:nonic31-known-binding}, then $D$ is the corresponding
global diamond divisor.
\end{lemma}
\begin{proof}
Let $B$ be the corresponding global divisor. The group
$J_1(31)(\mathbb Q)$ is finite by \cite[Theorem~3.1(2)]{DEHMZ}.
Since $X_1(31)$ has good reduction at $3$, reduction on rational torsion
is injective. The equality of reductions gives $[D-B]=0$ in
$J_1(31)(\mathbb Q)$. If $D\ne B$, a function with divisor $D-B$
would give a nonconstant map to $\mathbb P^1$ of degree at most $9$.
This contradicts $\operatorname{gon}_{\mathbb Q}X_1(31)=12$
\cite{DvH}. Hence $D=B$.
\end{proof}

\begin{corollary}\label{cor:nonic62-positive-rigidity}
Suppose that the finite class filter for a degree-$9$ point on $X_1(62)$
restricts the reduction of its image on $X_1(31)$ to the five places above.
Then this degree-$9$ point does not exist.
\end{corollary}
\begin{proof}
By Lemma~\ref{lem:nonic31-rigidity}, its image divisor is one of the
closed degree-$9$ divisors $B$. Consequently its residue field is the
full nonic field, rather than a proper subfield. Since $X_1(31)$ is a
fine moduli space, the underlying curve over this field is isomorphic
to the curve of the corresponding global diamond point. Its full torsion
is $\mathbb Z/31\mathbb Z$ by
Lemma~\ref{lem:nonic31-known-binding}. However, a point of order $62$
gives a nonzero rational point of order $2$, a contradiction.
\end{proof}
